% title:          Surjective monotone maps and continuity
% area:           limits-continuity
% methods:        construction, contradiction, monotonicity
% difficulty:     easy
% difficulty-ai:  easy
% status:
% review:         human
% --- history: all optional, leave EMPTY if unknown ---
% origin:         imc
% origin-date:    2006-07
% origin-number:  Day 1, Problem 1
% also-in:
% references:     Earliest known appearance (to the best of our knowledge): IMC 2006, Day 1, Problem 1 - https://www.imc-math.org.uk/imc2006/day1_solutions.pdf
% history:        ai
\begin{problem}
    Let \( f \colon \mathbb{R} \to \mathbb{R} \) be a function. Determine whether each of the
    following statements is true.
    \begin{enumerate}
        \item[(a)] If \( f \) is continuous and \( \operatorname{ran}\left(f\right) = \mathbb{R} \),
            then \( f \) is monotonic.
        \item[(b)] If \( f \) is monotonic and \( \operatorname{ran}\left(f\right) = \mathbb{R} \),
            then \( f \) is continuous.
        \item[(c)] If \( f \) is monotonic and continuous, then \( \operatorname{ran}\left(f\right)
            = \mathbb{R} \).
    \end{enumerate}
\end{problem}
\begin{solution}
    (a) False: the function \( x \mapsto x^2\left(x-1\right) \) is clearly continuous, and its range
    is \( \mathbb{R} \) by the intermediate value theorem, since it tends to \( -\infty \) at
    \( -\infty \) and to \( +\infty \) at \( +\infty \) (as \( x^2 \to +\infty \), while
    \( x-1 \to -\infty \) and \( x-1 \to +\infty \) respectively). But
    \[
        0^2\left(0-1\right) = 0, \quad \left(\frac{1}{2}\right)^2\left(\frac{1}{2}-1\right) =
        -\frac{1}{8} \quad \text{and} \quad 1^2\left(1-1\right) = 0,
    \]
    so that it is neither non-decreasing (its value at \( \frac{1}{2} \) is smaller than at \( 0 \))
    nor non-increasing (its value at \( 1 \) is bigger than at \( \frac{1}{2} \)), i.e., it is not
    monotonic.
    \\\\
    (b) True: without loss of generality, assume that \( f \) is increasing (otherwise, replace
    \( f \) by \( -f \), which is increasing, has range \( \mathbb{R} \) too, and is continuous if
    and only if \( f \) is). Denote by \( \mathcal{T}_{\mathbb{R}} \) the usual topology on
    \( \mathbb{R} \), and by \( \mathcal{T}_{\mathbb{R}}^{\ell} \) and
    \( \mathcal{T}_{\mathbb{R}}^{r} \) the topologies on \( \mathbb{R} \) generated by the intervals
    \( \left[x;y\right[ \) and \( \left]x;y\right] \) respectively, for \( x, y \in \mathbb{R} \)
    with \( x < y \).
    \\\\
    Now, assume by contradiction that \( f \) is not
    \( \mathcal{T}_{\mathbb{R}} \)-\( \mathcal{T}_{\mathbb{R}} \) continuous. Since
    \( \mathcal{T}_{\mathbb{R}}^{\ell} \cap \mathcal{T}_{\mathbb{R}}^{r} =
    \mathcal{T}_{\mathbb{R}} \) (a set that is open in both topologies contains, around each of its
    points \( x \), some \( \left]x-\delta;x\right] \) and some \( \left[x;x+\delta\right[ \) with
    \( \delta \in \mathbb{R}_{>0} \), hence \( \left]x-\delta;x+\delta\right[ \); the other
    inclusion is clear), this means that there exists \( U \in \mathcal{T}_{\mathbb{R}} \) with
    \( f^{-1}\left[U\right] \notin \mathcal{T}_{\mathbb{R}}^{\ell} \) or \( f^{-1}\left[U\right]
    \notin \mathcal{T}_{\mathbb{R}}^{r} \), i.e., \( f \) is not
    \( \mathcal{T}_{\mathbb{R}}^{\ell} \)-\( \mathcal{T}_{\mathbb{R}} \) continuous or \( f \) is
    not \( \mathcal{T}_{\mathbb{R}}^{r} \)-\( \mathcal{T}_{\mathbb{R}} \) continuous. Assume that
    \( f \) is not \( \mathcal{T}_{\mathbb{R}}^{\ell} \)-\( \mathcal{T}_{\mathbb{R}} \) continuous;
    the other case is the same, mutatis mutandis.
    \\\\
    Since \( \mathcal{T}_{\mathbb{R}}^{\ell} \) is first countable (each \( x \in \mathbb{R} \) has
    the neighbourhood basis \( \left\{ \left[x;x+\frac{1}{k}\right[ \,\middle|\, k \in
    \mathbb{N}_{\geq 1} \right\} \)), the
    \( \mathcal{T}_{\mathbb{R}}^{\ell} \)-\( \mathcal{T}_{\mathbb{R}} \) continuity of \( f \) is
    equivalent to its sequential continuity. Thus, there exist \( a \in \mathbb{R} \) and
    \( x_{-} \in \mathbb{R}^{\mathbb{N}} \) such that \( x_n
    \overset{n \to +\infty}{\longrightarrow} a \) in \( \mathcal{T}_{\mathbb{R}}^{\ell} \), but
    \( f\left(x_n\right) \) does not converge to \( f\left(a\right) \) in
    \( \mathcal{T}_{\mathbb{R}} \). Since \( \left[a;a+1\right[ \in
    \mathcal{T}_{\mathbb{R}}^{\ell} \) contains \( a \), there exists \( N \in \mathbb{N} \) such
    that \( x_n \geq a \) for all \( n \in \mathbb{N}_{\geq N} \), and, as
    \( \mathcal{T}_{\mathbb{R}} \subset \mathcal{T}_{\mathbb{R}}^{\ell} \), we have \( x_n
    \overset{n \to +\infty}{\longrightarrow} a \) in \( \mathcal{T}_{\mathbb{R}} \). As
    \( f\left(x_n\right) = f\left(a\right) \) whenever \( x_n = a \), and \( f\left(x_n\right) \)
    does not converge to \( f\left(a\right) \), infinitely many \( x_n \) with \( n \in
    \mathbb{N}_{\geq N} \) differ from \( a \), and their images do not converge to
    \( f\left(a\right) \) either (otherwise, all the images would). Replacing \( x_{-} \) by the
    subsequence of these terms, which still converges to \( a \) in both topologies, we may thus
    assume that \( x_n > a \) for all \( n \in \mathbb{N}_{\geq N} \). So, since \( f \) is
    increasing, \( \left\{ f\left(x_n\right) \,\middle|\, n \in \mathbb{N}_{\geq N} \right\} \subset
    \left[f\left(a\right);+\infty\right[ \), and it is standard (using \( x_n
    \overset{n \to +\infty}{\longrightarrow} a \) in \( \mathcal{T}_{\mathbb{R}} \)) that
    \( f\left(x_{-}\right) \) converges in \( \mathcal{T}_{\mathbb{R}} \) to
    \[
        \inf\left\{ f\left(x_n\right) \,\middle|\, n \in \mathbb{N}_{\geq N} \right\} \geq
        f\left(a\right).
    \]
    Moreover, as this limit is not \( f\left(a\right) \) by hypothesis, \( f\left(a\right) <
    \lim_{n \to +\infty} f\left(x_n\right) \), so let \( t \in
    \left]f\left(a\right);\lim_{n \to +\infty} f\left(x_n\right)\right[ \), which is non-empty.
    Since \( \operatorname{ran}\left(f\right) = \mathbb{R} \), there exists \( y \in \mathbb{R} \)
    with \( f\left(y\right) = t \). Now, for all \( n \in \mathbb{N}_{\geq N} \),
    \[
        f\left(a\right) < t = f\left(y\right) < \lim_{n' \to +\infty} f\left(x_{n'}\right) \leq
        f\left(x_n\right),
    \]
    where the last inequality holds since this limit is the infimum above, so that \( a < y <
    x_n \), because \( f \) is increasing (if \( y \leq a \), then \( f\left(y\right) \leq
    f\left(a\right) \), and if \( y \geq x_n \), then \( f\left(y\right) \geq f\left(x_n\right) \)).
    This is a contradiction to \( x_n \overset{n \to +\infty}{\longrightarrow} a \) in
    \( \mathcal{T}_{\mathbb{R}} \), as \( x_n-a > y-a > 0 \) for all \( n \in
    \mathbb{N}_{\geq N} \). Therefore, \( f \) is continuous.
    \\\\
    (c) False: the function \( f \colon \mathbb{R} \to \mathbb{R} \), \( x \mapsto
    \arctan\left(x\right) \) is monotonic and continuous, but
    \[
        \operatorname{ran}\left(f\right) = \left]-\frac{\pi}{2};\frac{\pi}{2}\right[ \subsetneq
        \mathbb{R}.
    \]
\end{solution}
